\documentclass[10pt,a4paper]{amsart}
\usepackage{amsmath,amssymb,amsthm,graphicx}
% Retained preamble for the Stacks Project chapter snapshot in this directory.
% The chapter's own preamble.tex is a copy of the shared Stacks Project preamble; the block
% below is sliced MECHANICALLY, line for line, from it. Nothing was added, removed or
% redefined; the wrapper class, the xr-hyper/hyperref wiring and the externaldocument list
% are omitted because the assembler supplies the document class.
%
% lines taken from preamble.tex: 3, 5, 9, 11, 12, 156, 157, 158, ... (51 lines total)
%
% package coverage restored from the chapter's own preamble (2026-10-13): xy, xy-2cell

\usepackage[all]{xy}
\xyoption{2cell}
\UseAllTwocells

\newenvironment{reference}{\comment}{\endcomment}
\newenvironment{slogan}{\comment}{\endcomment}
\newenvironment{history}{\comment}{\endcomment}
\theoremstyle{plain}
\newtheorem{theorem}[subsection]{Theorem}
\newtheorem{proposition}[subsection]{Proposition}
\newtheorem{lemma}[subsection]{Lemma}
\theoremstyle{definition}
\newtheorem{definition}[subsection]{Definition}
\newtheorem{example}[subsection]{Example}
\newtheorem{exercise}[subsection]{Exercise}
\newtheorem{situation}[subsection]{Situation}
\theoremstyle{remark}
\newtheorem{remark}[subsection]{Remark}
\newtheorem{remarks}[subsection]{Remarks}
\numberwithin{equation}{subsection}
\def\lim{\mathop{\mathrm{lim}}\nolimits}
\def\colim{\mathop{\mathrm{colim}}\nolimits}
\def\Spec{\mathop{\mathrm{Spec}}}
\def\Hom{\mathop{\mathrm{Hom}}\nolimits}
\def\Ext{\mathop{\mathrm{Ext}}\nolimits}
\def\SheafHom{\mathop{\mathcal{H}\!\mathit{om}}\nolimits}
\def\SheafExt{\mathop{\mathcal{E}\!\mathit{xt}}\nolimits}
\def\Sch{\mathit{Sch}}
\def\Mor{\mathop{\mathrm{Mor}}\nolimits}
\def\Ob{\mathop{\mathrm{Ob}}\nolimits}
\def\Sh{\mathop{\mathit{Sh}}\nolimits}
\def\NL{\mathop{N\!L}\nolimits}
\def\CH{\mathop{\mathrm{CH}}\nolimits}
\def\proetale{{pro\text{-}\acute{e}tale}}
\def\etale{{\acute{e}tale}}
\def\QCoh{\mathit{QCoh}}
\def\Ker{\mathop{\mathrm{Ker}}}
\def\Im{\mathop{\mathrm{Im}}}
\def\Coker{\mathop{\mathrm{Coker}}}
\def\Coim{\mathop{\mathrm{Coim}}}
\DeclareMathSymbol{\boxtimes}{\mathbin}{AMSa}{"02}
\def\QCohstack{\mathcal{QC}\!\mathit{oh}}
\def\Cohstack{\mathcal{C}\!\mathit{oh}}
\def\Spacesstack{\mathcal{S}\!\mathit{paces}}
\def\Quotfunctor{\mathrm{Quot}}
\def\Hilbfunctor{\mathrm{Hilb}}
\def\Curvesstack{\mathcal{C}\!\mathit{urves}}
\def\Polarizedstack{\mathcal{P}\!\mathit{olarized}}
\def\Complexesstack{\mathcal{C}\!\mathit{omplexes}}
\def\Pic{\mathop{\mathrm{Pic}}\nolimits}
\def\Picardstack{\mathcal{P}\!\mathit{ic}}
\def\Picardfunctor{\mathrm{Pic}}
\def\Deformationcategory{\mathcal{D}\!\mathit{ef}}
\usepackage{cleveref}
\crefname{originalproposition}{original proposition}{original propositions}
\Crefname{originalproposition}{Original proposition}{Original propositions}
\makeatletter
\newlabel{lemma-2-morphisms-presheaves}{{lemma-2-morphisms-presheaves}{source}{}{sourcepointer-1}{}}
\newlabel{lemma-2-morphisms-presheaves@cref}{{[originalproposition][0][]lemma-2-morphisms-presheaves}{[1][0][]source}{}{}{}}
\makeatother
\crefname{originalproposition}{original proposition}{original propositions}
\Crefname{originalproposition}{Original proposition}{Original propositions}
\makeatletter
\newlabel{lemma-characterize-quasi-coherent}{{lemma-characterize-quasi-coherent}{source}{}{sourcepointer-2}{}}
\newlabel{lemma-characterize-quasi-coherent@cref}{{[originalproposition][0][]lemma-characterize-quasi-coherent}{[1][0][]source}{}{}{}}
\makeatother
\crefname{originalproposition}{original proposition}{original propositions}
\Crefname{originalproposition}{Original proposition}{Original propositions}
\makeatletter
\newlabel{lemma-compare-quasi-coherent}{{lemma-compare-quasi-coherent}{source}{}{sourcepointer-3}{}}
\newlabel{lemma-compare-quasi-coherent@cref}{{[originalproposition][0][]lemma-compare-quasi-coherent}{[1][0][]source}{}{}{}}
\makeatother
\crefname{originalproposition}{original proposition}{original propositions}
\Crefname{originalproposition}{Original proposition}{Original propositions}
\makeatletter
\newlabel{lemma-stackification-quasi-coherent}{{lemma-stackification-quasi-coherent}{source}{}{sourcepointer-4}{}}
\newlabel{lemma-stackification-quasi-coherent@cref}{{[originalproposition][0][]lemma-stackification-quasi-coherent}{[1][0][]source}{}{}{}}
\makeatother
\crefname{originalproposition}{original proposition}{original propositions}
\Crefname{originalproposition}{Original proposition}{Original propositions}
\makeatletter
\newlabel{spaces-groupoids-definition-quotient-stack}{{044Q}{source}{}{sourcepointer-5}{}}
\newlabel{spaces-groupoids-definition-quotient-stack@cref}{{[originalproposition][0][]044Q}{[1][0][]source}{}{}{}}
\makeatother
\crefname{originalproposition}{original proposition}{original propositions}
\Crefname{originalproposition}{Original proposition}{Original propositions}
\makeatletter
\newlabel{spaces-groupoids-lemma-quotient-stack-2-arrow}{{044S}{source}{}{sourcepointer-6}{}}
\newlabel{spaces-groupoids-lemma-quotient-stack-2-arrow@cref}{{[originalproposition][0][]044S}{[1][0][]source}{}{}{}}
\makeatother
\crefname{originalproposition}{original proposition}{original propositions}
\Crefname{originalproposition}{Original proposition}{Original propositions}
\makeatletter
\newlabel{spaces-groupoids-lemma-quotient-stack-cocycle}{{044T}{source}{}{sourcepointer-7}{}}
\newlabel{spaces-groupoids-lemma-quotient-stack-cocycle@cref}{{[originalproposition][0][]044T}{[1][0][]source}{}{}{}}
\makeatother
\title{Stacks, Sheaves on Algebraic Stacks: proposition-quasi-coherent}
\author{The Stacks Project Authors}
\date{stacks-project:stacks-sheaves.tex:local-snapshot-20261012; GNU FDL 1.2 or later}
\begin{document}
\maketitle
\noindent\textit{Editorial note.} Statement span is lines 2350-2357 of sources/stacks-algebraic-stacks-sheaves/pinned-original-source.tex; proof span is lines 2359-2452, the complete author proof (92 body lines).\par\medskip
\noindent\textit{Editorial note.} Navigation stubs: 7 cited-but-not-included labels are supplied as non-load-bearing pointers with their pinned tags where the snapshot tag index has them.\par\medskip
\noindent\textit{Editorial note.} extra\_packages is empty on purpose: the retained preamble already loads xy with options and a second declaration clashes in the assembler wrapper.\par\medskip
\noindent\textit{Editorial note.} Text is retained verbatim, including the author's own typos and placeholders.\par\medskip
\noindent\textit{Editorial note.} Queue index 3 of queue\_algebraic-stacks-sheaves.json (pb=3913, stmt 2350-2357, proof 2359-2452).\par\medskip
\noindent\textit{Allowed standard prerequisites (editorial).}\begin{itemize}
\item Explicit prerequisite boundary. The cited inputs are prerequisites; the argument in the retained span is not replaced by any cited result.
\end{itemize}\medskip
\ifdefined\hypertarget\hypertarget{sourcepointer-1}{}\fi
\noindent\textit{Source reference.} Original proposition lemma-2-morphisms-presheaves: sources/stacks-algebraic-stacks-sheaves/pinned-original-source.tex, source lines 440--454. Cited by the retained statement or proof but not part of the collected spans; supplied only as a navigation pointer. The argument does not depend on its content.\par\medskip
\ifdefined\hypertarget\hypertarget{sourcepointer-2}{}\fi
\noindent\textit{Source reference.} Original proposition lemma-characterize-quasi-coherent: sources/stacks-algebraic-stacks-sheaves/pinned-original-source.tex, source lines 1798--1805. Cited by the retained statement or proof but not part of the collected spans; supplied only as a navigation pointer. The argument does not depend on its content.\par\medskip
\ifdefined\hypertarget\hypertarget{sourcepointer-3}{}\fi
\noindent\textit{Source reference.} Original proposition lemma-compare-quasi-coherent: sources/stacks-algebraic-stacks-sheaves/pinned-original-source.tex, source lines 2263--2274. Cited by the retained statement or proof but not part of the collected spans; supplied only as a navigation pointer. The argument does not depend on its content.\par\medskip
\ifdefined\hypertarget\hypertarget{sourcepointer-4}{}\fi
\noindent\textit{Source reference.} Original proposition lemma-stackification-quasi-coherent: sources/stacks-algebraic-stacks-sheaves/pinned-original-source.tex, source lines 2197--2207. Cited by the retained statement or proof but not part of the collected spans; supplied only as a navigation pointer. The argument does not depend on its content.\par\medskip
\ifdefined\hypertarget\hypertarget{sourcepointer-5}{}\fi
\noindent\textit{Source reference.} Original proposition 044Q: sources/stacks-algebraic-stacks-sheaves/pinned-original-source.tex, source lines 2452--2452. Cited by the retained statement or proof but not part of the collected spans; supplied only as a navigation pointer. The argument does not depend on its content.\par\medskip
\ifdefined\hypertarget\hypertarget{sourcepointer-6}{}\fi
\noindent\textit{Source reference.} Original proposition 044S: sources/stacks-algebraic-stacks-sheaves/pinned-original-source.tex, source lines 2452--2452. Cited by the retained statement or proof but not part of the collected spans; supplied only as a navigation pointer. The argument does not depend on its content.\par\medskip
\ifdefined\hypertarget\hypertarget{sourcepointer-7}{}\fi
\noindent\textit{Source reference.} Original proposition 044T: sources/stacks-algebraic-stacks-sheaves/pinned-original-source.tex, source lines 2452--2452. Cited by the retained statement or proof but not part of the collected spans; supplied only as a navigation pointer. The argument does not depend on its content.\par\medskip
\begin{proposition}
\label{proposition-quasi-coherent}
Let $(U, R, s, t, c)$ be a groupoid in algebraic spaces over $S$.
Let $\mathcal{X} = [U/R]$ be the quotient stack.
The category of quasi-coherent modules on $\mathcal{X}$
is equivalent to the category of quasi-coherent modules
on $(U, R, s, t, c)$.
\end{proposition}

\begin{proof}
We will construct quasi-inverse functors
$$
\QCoh(\mathcal{O}_\mathcal{X})
\longleftrightarrow
\QCoh(U, R, s, t, c).
$$
where $\QCoh(U, R, s, t, c)$ denotes the category of quasi-coherent modules
on the groupoid $(U, R, s, t, c)$.

\medskip\noindent
Let $\mathcal{F}$ be an object of $\QCoh(\mathcal{O}_\mathcal{X})$.
Denote $\mathcal{U}$, $\mathcal{R}$ the categories fibred in groupoids
corresponding to $U$ and $R$. Denote $x$ the (defining) object of $\mathcal{X}$
over $U$. Recall that we have a $2$-commutative diagram
$$
\xymatrix{
\mathcal{R} \ar[r]_s \ar[d]_t & \mathcal{U} \ar[d]^x \\
\mathcal{U} \ar[r]^x & \mathcal{X}
}
$$
See Groupoids in Spaces, Lemma
\ref{spaces-groupoids-lemma-quotient-stack-2-arrow}.
By Lemma \ref{lemma-2-morphisms-presheaves}
the $2$-arrow inherent in the diagram induces an isomorphism
$\alpha : t^*x^*\mathcal{F} \to s^*x^*\mathcal{F}$
which satisfies the cocycle condition over
$\mathcal{R} \times_{s, \mathcal{U}, t} \mathcal{R}$;
this is a consequence of Groupoids in Spaces, Lemma
\ref{spaces-groupoids-lemma-quotient-stack-cocycle}.
Thus if we set $\mathcal{G} = x^*\mathcal{F}|_{U_\etale}$
then the equivalence of categories in
Lemma \ref{lemma-compare-quasi-coherent}
(used several times compatibly with pullbacks)
gives an isomorphism
$\alpha : t_{small}^*\mathcal{G} \to s_{small}^*\mathcal{G}$
satisfying the cocycle condition on $R \times_{s, U, t} R$, i.e.,
$(\mathcal{G}, \alpha)$ is an object of $\QCoh(U, R, s, t, c)$.
The rule $\mathcal{F} \mapsto (\mathcal{G}, \alpha)$ is our
functor from left to right.

\medskip\noindent
Construction of the functor in the other direction.
Let $(\mathcal{G}, \alpha)$ be an object of $\QCoh(U, R, s, t, c)$.
According to Lemma \ref{lemma-stackification-quasi-coherent}
the stackification map $[U/_{\!p}R] \to [U/R]$ (see
Groupoids in Spaces, Definition
\ref{spaces-groupoids-definition-quotient-stack})
induces an equivalence of categories of quasi-coherent sheaves.
Thus it suffices to construct a quasi-coherent module $\mathcal{F}$
on $[U/_{\!p}R]$.

\medskip\noindent
Recall that an object $x = (T, u)$ of $[U/_{\!p}R]$
is given by a scheme $T$ and a morphism $u : T \to U$. A morphism
$(T, u) \to (T', u')$ is given by a pair $(f, r)$ where $f : T \to T'$
and $r : T \to R$ with $s \circ r = u$ and $t \circ r = u' \circ f$.
Let us call a {\it special morphism} any morphism of the form
$(f, e \circ u' \circ f) : (T, u' \circ f) \to (T', u')$.
The category of $(T, u)$ with special morphisms is just the
category of schemes over $U$.

\medskip\noindent
With this notation in place, given an object $(T, u)$ of $[U/_{\!p}R]$, we set
$$
\mathcal{F}(T, u) : = \Gamma(T, u_{small}^*\mathcal{G}).
$$
Given a morphism $(f, r) : (T, u) \to (T', u')$ we get a map
\begin{align*}
\mathcal{F}(T', u') & = \Gamma(T', (u')_{small}^*\mathcal{G}) \\
& \to \Gamma(T, f_{small}^*(u')_{small}^*\mathcal{G}) =
\Gamma(T, (u' \circ f)_{small}^*\mathcal{G}) \\
& = \Gamma(T, (t \circ r)_{small}^*\mathcal{G}) =
\Gamma(T, r_{small}^*t_{small}^*\mathcal{G}) \\
& \to \Gamma(T, r_{small}^*s_{small}^*\mathcal{G}) =
\Gamma(T, (s \circ r)_{small}^*\mathcal{G}) \\
& = \Gamma(T, u_{small}^*\mathcal{G}) \\
& = \mathcal{F}(T, u)
\end{align*}
where the first arrow is pullback along $f$ and the second arrow is
$\alpha$. Note that if $(f, r)$ is a special morphism, then this
map is just pullback along $f$ as $e_{small}^*\alpha = \text{id}$ by
the axioms of a sheaf of quasi-coherent modules on a groupoid.
The cocycle condition implies that $\mathcal{F}$ is a presheaf
of modules (details omitted). We see that the restriction of $\mathcal{F}$
to $(\Sch/T)_{fppf}$ is quasi-coherent by the simple description of the
restriction maps of $\mathcal{F}$ in case of a special morphism.
Hence $\mathcal{F}$ is a sheaf on $[U/_{\!p}R]$ and quasi-coherent
(Lemma \ref{lemma-characterize-quasi-coherent}).

\medskip\noindent
We omit the verification that the functors constructed above are
quasi-inverse to each other.
\end{proof}

\end{document}
